\documentclass{beamer} % Classe de document pour les présentations modernes
\usepackage[utf8]{inputenc} % Encodage des caractères
\usepackage{graphicx} % Pour inclure des images

% Thème de la présentation (vous pouvez le changer)
\usetheme{Madrid} % Thème moderne
\usecolortheme{seahorse} % Palette de couleurs

% Informations de la présentation
\title{Titre de la Présentation}
\subtitle{Sous-titre (facultatif)}
\author{Prénom Nom}
\institute{Nom de l'Établissement}
\date{\today}

\begin{document}

% Page de titre
\begin{frame}
    \titlepage
\end{frame}

% Table des matières
\begin{frame}
    \frametitle{Plan de la Présentation}
    \tableofcontents % Génère automatiquement une table des matières
\end{frame}

% Section 1 : Introduction
\section{Introduction}
\begin{frame}
    \frametitle{Introduction}
    \framesubtitle{Contexte et Objectifs}
    \begin{itemize}
        \item<1-> \textbf{Contexte} : Description du contexte.
        \item<2-> \textbf{Objectifs} : Objectifs de la présentation.
        \item<3-> \textbf{Plan} : Vue d'ensemble du plan.
    \end{itemize}
\end{frame}

% Section 2 : Méthodologie
\section{Méthodologie}
\begin{frame}
    \frametitle{Méthodologie}
    \begin{enumerate}
        \item Étape 1 : Description de l'étape 1.
        \item Étape 2 : Description de l'étape 2.
        \item Étape 3 : Description de l'étape 3.
    \end{enumerate}
\end{frame}

% Section 3 : Résultats
\section{Résultats}
\begin{frame}
    \frametitle{Résultats}
    \begin{block}{Résultats clés}
        \begin{itemize}
            \item Résultat 1 : Description.
            \item Résultat 2 : Description.
        \end{itemize}
    \end{block}
\end{frame}

% Section 4 : Conclusion
\section{Conclusion}
\begin{frame}
    \frametitle{Conclusion}
    \begin{itemize}
        \item Synthèse des points clés.
        \item Perspectives futures.
    \end{itemize}
\end{frame}

\end{document}

